\chapter{Cadre g\'en\'eral du projet}
\finPageTitre

\section{Introduction}
Ce premier chapitre est consacr\'e \`a la pr\'esentation du cadre g\'en\'eral du projet de fin d'ann\'ee. Nous introduisons tout d'abord l'organisme d'accueil, ses activit\'es et ses p\^oles d'expertise. Nous menons ensuite une analyse m\'ethodique de l'existant afin de cerner le contexte technique et op\'erationnel. Enfin, nous d\'elimitons le p\'erim\`etre pr\'ecis de la mission confi\'ee, ses objectifs, ses contraintes ainsi que la planification pr\'evisionnelle illustr\'ee par le diagramme de Gantt.

\section{Pr\'esentation de l'organisme d'accueil}
\subsection{Pr\'esentation g\'en\'erale et secteur d'activit\'e}
L'organisme d'accueil, \textbf{[Nom de l'organisme / Entreprise]}, est une structure sp\'ecialis\'ee dans le domaine de \textbf{[secteur d'activit\'e : ex. le d\'eveloppement logiciel, l'int\'egration de syst\`emes, la transformation digitale]}.

Fond\'ee en \textbf{[Ann\'ee]}, l'organisation propose des solutions innovantes adapt\'ees aux besoins de ses clients et partenaires. Ses activit\'es couvrent le conseil en technologies, la conception d'architectures logicielles et le d\'eploiement de solutions m\'etier.

\subsection{Positionnement et \'ecosyst\`eme}
L'organisme collabore avec divers partenaires institutionnels et industriels \textbf{[citer des r\'ef\'erences : ex. entreprises partenaires, organismes publics, clients cl\'es]}. Situ\'e \`a \textbf{[Ville / R\'egion]}, il dispose d'un positionnement strat\'egique lui permettant d'intervenir sur des projets \`a forte valeur ajout\'ee.

\subsection{Domaines d'expertise}
Les principales comp\'etences et activit\'es de l'organisme s'articulent autour de plusieurs axes compl\'ementaires pr\'esent\'es dans le Tableau~\ref{tab:expertises}.

\begin{table}[H]
\centering
\begin{tabularx}{\textwidth}{X X X}
\toprule
\textbf{Ing\'enierie Logicielle} & \textbf{Conseil \& Expertise} & \textbf{Innovation \& R\&D} \\
\midrule
\textbullet~Conception d'architectures \par
\textbullet~D\'eveloppement Web \& Mobile \par
\textbullet~Int\'egration de syst\`emes & 
\textbullet~Audit et s\'ecurit\'e \par
\textbullet~Gestion de projets IT \par
\textbullet~Accompagnement digital & 
\textbullet~Veille technologique \par
\textbullet~D\'eveloppement de PoC \par
\textbullet~Intelligence Artificielle \\
\bottomrule
\end{tabularx}
\caption{Domaines d'expertise de l'organisme d'accueil}
\label{tab:expertises}
\end{table}

\subsection{Solutions et produits principaux}
Dans le cadre de ses prestations, l'organisme con\c{c}oit et maintient diff\'erentes solutions logicielles :
\begin{itemize}
    \item \textbf{[Solution / Plateforme 1]} : [Br\`eve description de la solution principale ou du produit associ\'e au projet] ;
    \item \textbf{[Solution / Plateforme 2]} : [Br\`eve description d'autres produits d\'evelopp\'es au sein de l'entreprise].
\end{itemize}

\section{Analyse de l'existant}
\subsection{Contexte et p\'erim\`etre du syst\`eme existant}
L'intervention r\'ealis\'ee porte sur \textbf{[Nom du syst\`eme / de l'application existante]}. Ce syst\`eme a pour finalit\'e de \textbf{[d\'ecrire la mission op\'erationnelle du syst\`eme existant]}.

\subsection{Architecture technique initiale}
Le Tableau~\ref{tab:stack-existant} r\'esume la composition de la pile technologique initiale du syst\`eme.

\begin{table}[H]
\centering
\begin{tabularx}{\textwidth}{p{3.8cm} X p{3.8cm}}
\toprule
\textbf{Composant / Couche} & \textbf{Technologies Utilis\'ees} & \textbf{R\^ole Op\'erationnel} \\
\midrule
\textbf{Couche Pr\'esentation} & [Technologies Front : ex. HTML/CSS, Framework UI] & Interface utilisateur et interactions \\
\textbf{Couche M\'etier}       & [Technologies Back : ex. Langage, Framework API]  & Logique m\'etier et traitement des flux \\
\textbf{Couche Persistance}    & [Syst\`eme de BDD : ex. SQL / NoSQL]               & Stockage et gestion des donn\'ees \\
\textbf{Infrastructure}        & [Outils : ex. Serveur web, Conteneurs, Outils CI]   & H\'ebergement et ex\'ecution \\
\bottomrule
\end{tabularx}
\caption{Architecture technique du syst\`eme existant}
\label{tab:stack-existant}
\end{table}

\subsection{Acteurs et flux fonctionnels}
Le fonctionnement de la plateforme implique diff\'erents r\^oles d'utilisateurs :
\begin{itemize}
    \item \textbf{[Acteur 1 : ex. Administrateur]} : [R\^ole, responsabilit\'es et droits d'acc\`es] ;
    \item \textbf{[Acteur 2 : ex. Gestionnaire / Utilisateur m\'etier]} : [Op\'erations principales et t\^aches quotidiennes] ;
    \item \textbf{[Acteur 3 : ex. Client / Utilisateur final]} : [Consultation des informations et interactions avec le syst\`eme].
\end{itemize}

\begin{figure}[H]
\centering
\IfFileExists{figures/ch1/diagramme_fonctionnel.png}{%
    \includegraphics[width=0.85\textwidth]{figures/ch1/diagramme_fonctionnel.png}%
}{%
    \fbox{\parbox[c][5cm][c]{0.85\textwidth}{\centering \textit{[Ins\'erer ici le diagramme de cas d'utilisation ou d'architecture fonctionnelle de l'existant]}}}%
}
\caption{Diagramme fonctionnel des flux du syst\`eme}
\label{fig:diag-fonctionnel}
\end{figure}

\section{Probl\'ematique et identification des besoins}
L'analyse de l'existant et les \'echanges avec l'\'equipe encadrante ont permis de d\'egager plusieurs limitations :
\begin{enumerate}
    \item \textbf{Limitations fonctionnelles} : [D\'ecrire les manques fonctionnels ou processus manuels \`a automatiser] ;
    \item \textbf{Contraintes techniques} : [D\'ecrire les contraintes d'architecture, de scalabilit\'e ou de performance] ;
    \item \textbf{Besoins d'am\'elioration} : [D\'ecrire les attentes prioritaires des utilisateurs et de l'organisme].
\end{enumerate}

\section{Expos\'e de la mission confi\'ee}
\subsection{Objectifs g\'en\'eraux}
Pour r\'epondre aux besoins identifi\'es, la mission confi\'ee dans le cadre de ce stage consiste \`a \textbf{[d\'efinir l'objectif majeur du projet : ex. concevoir et d\'evelopper un module m\'etier complet, refondre l'architecture applicative ou cr\'eer une nouvelle plateforme]}.

\subsection{P\'erim\`etre du projet}
\begin{itemize}
    \item \textbf{T\^aches et fonctionnalit\'es incluses} :
    \begin{itemize}
        \item [T\^ache 1 : ex. Analyse d\'etaill\'ee des besoins et r\'edaction des sp\'ecifications] ;
        \item [T\^ache 2 : ex. Conception de l'architecture logicielle et de la base de donn\'ees] ;
        \item [T\^ache 3 : ex. D\'eveloppement et int\'egration des composants Front-End et Back-End] ;
        \item [T\^ache 4 : ex. Mise en \oe uvre des tests de validation et d\'eploiement].
    \end{itemize}
    \item \textbf{Fonctionnalit\'es hors p\'erim\`etre} :
    \begin{itemize}
        \item [Exclusion 1 : Fonctionnalit\'es secondaires report\'ees aux versions futures] ;
        \item [Exclusion 2 : Int\'egrations mat\'erielles ou d\'eploiements externes hors port\'ee].
    \end{itemize}
\end{itemize}

\subsection{Contraintes et crit\`eres d'\'evaluation}
Le projet est soumis aux contraintes suivantes :
\begin{itemize}
    \item \textbf{D\'elai d'ex\'ecution} : Respect de la dur\'ee impartie pour le stage ([nombre] semaines) ;
    \item \textbf{Qualit\'e et s\'ecurit\'e} : Respect des standards de codage, de s\'ecurit\'e des donn\'ees et de modularit\'e ;
    \item \textbf{Maintenabilit\'e} : Production d'un code document\'e, testable et ais\'ement r\'eutilisable.
\end{itemize}

\section{Planification temporelle et diagramme de Gantt}
Afin de structurer le d\'eroulement du projet et d'assurer le respect des d\'elais, un planning pr\'evisionnel a \'et\'e \'etabli. La Figure~\ref{fig:gantt} pr\'esente la r\'epartition temporelle des principales phases du projet.

\begin{figure}[H]
\centering
\IfFileExists{figures/ch1/gantt.png}{%
    \includegraphics[width=0.92\textwidth]{figures/ch1/gantt.png}%
}{%
    \begin{ganttchart}[
        hgrid,
        vgrid,
        x unit=0.88cm,
        y unit chart=0.6cm,
        title height=1,
        bar/.append style={fill=grisfonce},
        bar height=0.6
    ]{1}{10}
    \gantttitle{Semaines}{10}\\
    \gantttitlelist{1,...,10}{1}\\
    \ganttbar{Analyse de l'existant \& Sp\'ecifications}{1}{2}\\
    \ganttbar{\'Etat de l'art \& Choix techniques}{2}{4}\\
    \ganttbar{Conception de l'architecture}{4}{6}\\
    \ganttbar{D\'eveloppement \& Int\'egration}{5}{8}\\
    \ganttbar{Tests, Validation \& D\'eploiement}{8}{9}\\
    \ganttbar{R\'edaction du rapport \& Cl\^oture}{9}{10}
    \end{ganttchart}%
}
\caption{Diagramme de Gantt pr\'evisionnel du d\'eroulement du projet}
\label{fig:gantt}
\end{figure}

\sectionTransition{%
Ce premier chapitre a pr\'esent\'e l'organisme d'accueil, analys\'e l'existant et d\'efini les objectifs et la planification du projet. Le deuxi\`eme chapitre est consacr\'e \`a l'\'etude de l'\'etat de l'art, \`a la d\'emarche m\'ethodologique et \`a la justification des choix technologiques et architecturaux.
}
